\section{Finite-prior failure at every exponent at least one}
\label{sec:relu-native-covariance}

The lower bound extends to every canonical exponent $\gamma\geq1$.
The comparison distribution below retains the prior's full covariance.
It only translates the Gaussian means.  Signed pairs of ReLU
populations transport their conditional mean difference exactly, even
when the common Gaussian background variance is large.

\begin{lemma}[A native-covariance Boolean memorizer]
\label{lem:relu-native-memorizer}
Let $H\geq2$, $L=H+1$, $1\leq D\leq d$, and
$r:\{-1,1\}^D\to[-1,1]$ satisfy $\sum_xr(x)^2\leq n$.
Put
\begin{gather*}
 M=(D+1)2^D,\qquad Q_1=4d(D+1)^2,\qquad
 Q_2=40(D+1)^2n,\\
 R^2=(Q_1Q_2\,6^{L-2})^{1/L},\qquad
 \beta=R/\sqrt{Q_1},\qquad
 e_n=\frac{\sqrt{20}(D+1)\sqrt n}{\beta}
                                      e^{-\beta^2/16}.
\end{gather*}
For every $N\geq M$ and every $\sigma^2$ with
$\tau^2=N\sigma^2\leq1$, there is a Gaussian comparison law $\nu$
with exactly the covariance of the prior $\mu_{N,\sigma}$ such that
\begin{gather}
 2\sigma^2\operatorname{KL}(\nu\Vert\mu_{N,\sigma})\leq LR^2,
 \label{eq:relu-native-entropy}\\
 |\mathbb E_\nu z_w(x)-r(x_{[D]})|\leq e_n,\qquad
 \operatorname{Var}_\nu z_w(x)
                 \leq\frac{12H(R^2+1)^H}{N}
 \label{eq:relu-native-accuracy}
\end{gather}
whenever the designated $D$ coordinates $x_{[D]}$ are Boolean and
all other coordinates lie in $[-1,1]$.
\end{lemma}

\begin{proof}
We first construct features whose hinge arguments stay away from zero
at every Boolean input.  For a parity on $j\geq1$ coordinates, let
$t$ be their sum, $t_i=-j+2i$, and $y_i=y_0(-1)^i$ its values.
The identity
\begin{equation}
 \chi(t)=2y_0\relu(t+j+1/2)
             -4\sum_{i=0}^{j-1}y_i\relu(t-t_i-1/2)
 \label{eq:relu-margin-parity}
\end{equation}
holds at $t=t_0,\ldots,t_j$.
For example it follows by shifting the triangular interpolant by
one half and multiplying by two, or directly by summing the alternating
slopes.  Every hinge argument has absolute value at least $1/2$ on
these points.  The constant parity is $2\relu(1/2)$.
Each parity uses at most $D+1$ features, each augmented row has norm
at most $2(D+1)$, and its coefficient squared sum is at most
$4+16j\leq20(D+1)$.

Apply this to the normalized Fourier expansion of $r$.  Keeping
separate copies of all features, and adding constant features with
coefficient zero if necessary, gives exactly $M$ affine forms $\ell_j$
and coefficients $c_j$ such that
\[
 r(x_{[D]})=\sum_{j=1}^M c_j\relu(\ell_j(x)),\qquad
 |\ell_j(x)|\geq1/2,\qquad
 T_c:=\sum_jc_j^2\leq20(D+1)n/2^D.
\]
The features ignore the other input coordinates.
Partition the first hidden layer into $M$ balanced nonempty groups,
with sizes $k_j\geq N/(2M)$.  Set each row mean in group $j$ to
$\beta\ell_j$.  Let $F_1=\sum_jc_j\overline h^1_j$, where the bar
denotes the empirical average in that group.

Split each later hidden layer into two groups of sizes $k_+,k_-$,
differing by at most one.  For the second layer use incoming row
means $+a_2c_j/k_j$ in its positive group and $-a_2c_j/k_j$ in its
negative group, where $a_2=R/\sqrt{Q_2}$.
For each layer $3\leq\ell\leq H$, put
$a_\ell=R/\sqrt6$.  A positive row has incoming means
$+a_\ell/k_+$ on the previous positive group and
$-a_\ell/k_-$ on its negative group; a negative row has their negatives.
Finally set $a_L=R/\sqrt6$ and readout means
$+a_L/k_+$ and $-a_L/k_-$.
All covariances remain exactly those of the original independent
Gaussian prior.

The first mean-layer energy is at most $Q_1\beta^2=R^2$.
The second is at most
$2MT_ca_2^2\leq Q_2a_2^2=R^2$.
Every subsequent internal mean-layer energy and the readout mean
energy is at most $6a_\ell^2=R^2$, because
$N(1/k_++1/k_-)\leq6$.
Moreover
\[
 \beta\prod_{\ell=2}^La_\ell
 =\frac{R^L}{\sqrt{Q_1Q_2\,6^{L-2}}}=1.
\]
The mean weights therefore have total energy at most $LR^2$, proving
\eqref{eq:relu-native-entropy} by the exact Gaussian shift formula.
In particular every internal mean matrix has operator norm at most $R$.

For any $u$ and any centered symmetric scalar random variable $Z$,
\[
 \mathbb E\relu(u+Z)-\mathbb E\relu(-u+Z)=u.
\]
This follows from $\relu(t)-\relu(-t)=t$ and symmetry.
Conditionally on the preceding layers, all fields in a paired layer
have the same Gaussian variance
$q=\sigma^2\|h^{\ell-1}(x)\|_2^2$ and means $\pm u$.
If $D_\ell$ is the positive-group average minus the negative-group
average, it follows that
\[
 \mathbb E[D_2\mid h^1]=a_2F_1,\qquad
 \mathbb E[D_\ell\mid h^{\ell-1}]=a_\ell D_{\ell-1}
                    \quad(3\leq\ell\leq H).
\]
The readout conditional mean is $a_LD_H$.

At the first layer every Gaussian field has variance at most two.
For $s\leq\sqrt2$ and $|u|\geq\beta/2$,
\[
 0\leq\mathbb E\relu(u+sZ)-\relu(u)
      \leq s(2\pi)^{-1/2}e^{-u^2/(2s^2)}
      \leq e^{-\beta^2/16},
 \qquad Z\sim N(0,1).
\]
Thus $|\mathbb E F_1-\beta r|\leq\sum_j|c_j|e^{-\beta^2/16}$.
Since $\sum_j|c_j|\leq\sqrt{MT_c}\leq\sqrt{20}(D+1)\sqrt n$,
the conditional identities and product of gains prove the mean bound
in~\eqref{eq:relu-native-accuracy}.

We check the finite-width variance explicitly.  For any scalar
one-Lipschitz function $f$ and random variable $Z$,
$\operatorname{Var}(f(Z))\leq\operatorname{Var}(Z)$, by using an
independent copy in the identity for variance.
Consequently
\[
 \operatorname{Var}F_1\leq2\sum_jc_j^2/k_j
             \leq4MT_c/N\leq2Q_2/N.
\]
Let $S_\ell=\|h^\ell(x)\|_2^2/N$.
The first mean-layer norm and the conditional Gaussian covariance give
\[
 \mathbb E S_1\leq2(R^2+1),\qquad
 \mathbb E S_\ell\leq(R^2+1)\mathbb E S_{\ell-1}
                 \leq2(R^2+1)^\ell.
\]
Here the cross term vanishes because the Gaussian perturbation is
centered, and $\relu(t)^2\leq t^2$.
Conditional independence of the rows also gives
\[
 \operatorname{Var}(D_\ell\mid h^{\ell-1})
 \leq(1/k_++1/k_-)\tau^2S_{\ell-1}
 \leq6S_{\ell-1}/N.
\]
Its expectation is at most $12(R^2+1)^{\ell-1}/N$.
The native readout noise contributes at most
$\mathbb E S_H/N\leq2(R^2+1)^H/N$.

Iterating the law of total variance, the first-layer contribution to
output variance is at most
\[
 \beta^{-2}(2Q_2/N)
 =\frac{2Q_1Q_2}{R^2N}
 =\frac{2R^{2H}}{6^{L-2}N}
 \leq\frac{2(R^2+1)^H}{N}.
\]
Each of the $H-1$ paired hidden layers contributes at most
$12(R^2+1)^H/N$ after the later gains, since each remaining squared
gain is at most $R^2+1$.  Adding the readout contribution proves the
stated upper bound $12H(R^2+1)^H/N$.
\end{proof}

\begin{theorem}[All-temperature failure for every finite $\gamma\geq1$]
\label{thm:relu-native-all-gamma}
Use the laws and sample event of Theorem~\ref{thm:lower}.
For every sufficiently large $n$, there is a deterministic polynomial
width threshold $N_{\rm nat}(n)$ such that, on that same sample event,
simultaneously for every $N\geq N_{\rm nat}(n)$, every $\gamma\geq1$,
and every $\lambda>0$, all the raw and clipped prediction risks stated
in Theorem~\ref{thm:lower} are at least $p+1/100$.
One may take, for $H\geq2$,
\[
 N_{\rm nat}(n)=\max\left\{m,(D+1)2^D,
                  \left\lceil12Hn(R^2+1)^H\right\rceil\right\},
 \qquad D=\log_2n+1,
\]
with $R$ as in Lemma~\ref{lem:relu-native-memorizer} and $d=D+1$.
Thus $N_{\rm nat}=O_H(m+n^{2-1/L}\operatorname{polylog}n)$.
For $H=1$, the earlier canonical finite-prior theorem already covers
all these exponents.
\end{theorem}
\begin{proof}
For the low-temperature branch set
$\tau_n=n^{-1/L}d^{-4}$ as before.
Use the empirical noisy means as $r$, assigning zero to all
negative-tag Boolean points.  Their squared counting norm is at most
$n$, so Lemma~\ref{lem:relu-native-memorizer} applies to the same
memorization problem.  Its normalized mean-energy cost is
\[
 LR^2=O_H(n^{1/L}d^{5/L}),\qquad \tau_nLR^2=o(1).
\]
Furthermore
$\beta^2=R^2/[4d(D+1)^2]$ grows as a positive power of $n$ up to
logarithms, so $e_n=o(n^{-2})$.
The chosen width makes the comparison output variance at most $1/n$.
At every sample point its expected squared loss exceeds the
memorizer loss by at most $4e_n+e_n^2+1/n$, and this increment is
at most $2/n$ for sufficiently large $n$.
Posterior variational optimality gives
\[
 \int\widehat R_S(z_w)d\Pi
 \leq\widehat R_S(r)+2/n+\lambda LR^2.
\]
This inequality is uniform in $N\geq N_{\rm nat}$ and $\gamma\geq1$,
because the comparison covariance is exactly the native one.
For $\lambda\leq\tau_n$, the posterior noisy residual is at most
$q+o(1)$ after subtracting the within-atom residual.
The occupancy event and the same Jensen/reverse-triangle argument
used in Theorem~\ref{thm:lower} give excess at least $1/100$ for all
its stated prediction conventions.

For completeness, the high-temperature argument is also uniform in
all canonical exponents.  Normalize the matrices by
\[
 A_1=\sqrt{d/N}\,W^1,\qquad A_\ell=W^\ell\ (2\leq\ell\leq H),
 \qquad A_L=\sqrt N\,v^{\mathsf T}.
\]
Every coordinate has variance $\sigma^2=N^{-\gamma}$.
With $T=\max_\ell\|A_\ell\|_{\rm op}$ and
$M_d=1+\sqrt{d+1}$, Gaussian operator-norm concentration gives
\[
 \mu\{T\geq M_d+t\}\leq L e^{-t^2/(2\sigma^2)}.
\]
The prior expected empirical risk is at most three, so
$Z\geq e^{-3/(2\lambda\sigma^2)}$.
The resulting posterior tail, integrated by comparison with a
mean-one exponential variable, implies
\[
 \left(\int T^L d\Pi\right)^{1/L}
 \leq M_d+\sqrt{3/\lambda+2\log L}+2\sqrt L.
\]
The input-coordinate Lipschitz constant is at most $d^{-1/2}T^L$.
For $\lambda\geq\tau_n$, the inner steep pair therefore has raw-mean
half-difference at most
\[
 \delta d^{-1/2}
 [M_d+\sqrt{3/\tau_n+2\log L}+2\sqrt L]^L
 =O_H(n^{1/2-Ha-Lb}\operatorname{polylog}n)=o(1).
\]
The same estimate applies to its clipping and to the mean of clipped
outputs; Jensen bounds the corresponding Gibbs risks from below.
The inner-pair argument again gives excess at least $q/8>1/100$.
Both temperature ranges use exactly the same law and sample event,
completing the proof.
\end{proof}

\begin{corollary}[Localized noise and fixed comparator width]
\label{cor:relu-native-variants}
The localized-noise conclusion of Corollary~\ref{cor:localized-noise}
also holds simultaneously for every finite $\gamma\geq1$, every
$\lambda>0$, and every $N\geq N_{\rm nat}(n)$, with the same excess
$c(p)=p^2e^{-p}/256$ for each fixed $0<p\leq9/10$.

Likewise, the fixed-width-four conclusion of
Corollary~\ref{cor:fixed-width} holds throughout the same finite-prior
range whenever $b>1/(2L)$, with the width-four teacher and with $m=4$
in $N_{\rm nat}(n)$.  The sufficiently large sample thresholds may
depend on the fixed depth and on $p$.
\end{corollary}
\begin{proof}
For localized noise, apply the same native-covariance memorizer to the
empirical noisy means, setting its target to zero on every negative-tag
point.  The extra comparison loss is at most $2/n$, and the normalized
entropy cost is still $LR^2$.  Thus at $\lambda\leq\tau_n$ the noisy
residual differs by $o(1)$ from the inner-pair sample mass used in
Corollary~\ref{cor:localized-noise}.  Its occupancy estimates and
reverse-triangle proof give exactly the stated constant $c(p)$.
The all-exponent high-temperature estimate above has the same vanishing
inner-pair half-difference, so that corollary's second branch is also
unchanged.

For a width-four teacher, the memorization problem and its comparison
law are again unchanged.  The steep separation is now
$\delta=2B^{-L}$, and the high-temperature estimate becomes
$O_H(n^{1/2-Lb}\operatorname{polylog}n)=o(1)$ under $Lb>1/2$.
This proves the fixed-width assertion without an assumption $a>0$.
All comparisons are uniform in the displayed widths, exponents, and
temperatures, and use the same sample events as the respective original
constructions.
\end{proof}
